\begin{thebibliography}{65}
\bibitem{ref1} Scott Aaronson and Avi Wigderson. Algebrization: A new barrier in complexity theory. ACM Transactions on Computation Theory, 1(1):2, 2009.
\bibitem{ref2} Miklós Ajtai. $\Sigma^1_1$-formulae on finite structures. Annals of pure and applied logic, 24(1):1-48, 1983.
\bibitem{ref3} Miklós Ajtai and Michael Ben-Or. A Theorem on Probabilistic Constant Depth Computations. In Proceedings of the Symposium on Theory of Computing (STOC), pages 471-474, 1984.
\bibitem{ref4} Eric Allender. The new complexity landscape around circuit minimization. In Alberto Leporati, Carlos Martín-Vide, Dana Shapira, and Claudio Zandron, editors, Language and Automata Theory and Applications (LATA), volume 12038 of Lecture Notes in Computer Science, pages 3-16. Springer, 2020.
\bibitem{ref5} Eric Allender and Bireswar Das. Zero knowledge and circuit minimization. In International Symposium on Mathematical Foundations of Computer Science (MFCS), pages 25-32, 2014.
\bibitem{ref6} Eric Allender, Joshua A. Grochow, Dieter van Melkebeek, Cristopher Moore, and Andrew Morgan. Minimum circuit size, graph isomorphism, and related problems. SIAM Journal on Computing, 47(4):1339-1372, 2018.
\bibitem{ref7} Eric Allender, Lisa Hellerstein, Paul McCabe, Toniann Pitassi, and Michael E. Saks. Minimizing disjunctive normal form formulas and AC$^0$ circuits given a truth table. SIAM Journal on Computing, 38(1):63-84, 2008.
\bibitem{ref8} Eric Allender and Shuichi Hirahara. New insights on the (non-)hardness of circuit minimization and related problems. In International Symposium on Mathematical Foundations of Computer Science (MFCS), pages 54:1-54:14, 2017.
\bibitem{ref9} Eric Allender, Dhiraj Holden, and Valentine Kabanets. The minimum oracle circuit size problem. Computational Complexity, 26(2):469-496, 2017.
\bibitem{ref10} Eric Allender, Rahul Ilango, and Neekon Vafa. The non-hardness of approximating circuit size. In International Computer Science Symposium in Russia (CSR), pages 13-24, 2019.
\bibitem{ref11} Eric Allender and Michal Koucky. Amplifying lower bounds by means of self-reducibility. Journal of the ACM, 57(3):1-36, 2010.
\bibitem{ref12} Eric Allender, Michal Koucky, Detlef Ronneburger, and Sambuddha Roy. The pervasive reach of resource-bounded Kolmogorov complexity in computational complexity theory. Journal of Computer and System Sciences, 77(1):14-40, 2011.
\bibitem{ref13} Noga Alon and Joel H Spencer. The probabilistic method. John Wiley \& Sons, 2016.
\bibitem{ref14} Theodore Baker, John Gill, and Robert Solovay. Relativizations of the p=?np question. SIAM Journal on computing, 4(4):431-442, 1975.
\bibitem{ref15} Boaz Barak, Mark Braverman, Xi Chen, and Anup Rao. How to compress interactive communication. SIAM Journal on Computing, 42(3):1327-1363, 2013.
\bibitem{ref16} Avrim L. Blum and Ronald L. Rivest. Training a 3-node neural network is np-complete. Neural Networks, 5(1):117-127, 1992.
\bibitem{ref17} Joan Boyar, Philip Matthews, and René Peralta. On the shortest linear straight-line program for computing linear forms. In International Symposium on Mathematical Foundations of Computer Science (MFCS), pages 168-179, 2008.
\bibitem{ref18} Marco L. Carmosino, Russell Impagliazzo, Valentine Kabanets, and Antonina Kolokolova. Learning algorithms from natural proofs. In Conference on Computational Complexity (CCC), pages 10:1-10:24, 2016.
\bibitem{ref19} Sebastian Lukas Arne Czort. The complexity of minimizing disjunctive normal form formulas. Master’s thesis, University of Aarhus, 1999.
\bibitem{ref20} Uriel Feige and Joe Kilian. Zero knowledge and the chromatic number. Journal of Computer and System Sciences, 57(2):187-199, 1998.
\bibitem{ref21} Vitaly Feldman. Hardness of approximate two-level logic minimization and PAC learning with membership queries. In Symposium on Theory of Computing (STOC), pages 363-372, 2006.
\bibitem{ref22} Alexander Golovnev, Rahul Ilango, Russell Impagliazzo, Valentine Kabanets, Antonina Kolokolova, and Avishay Tal. ac$^0$[p] lower bounds against mcsp via the coin problem. In Colloquium on Automata, Languages, and Programming (ICALP), volume 132, page 66, 2019.
\bibitem{ref23} Thomas R. Hancock, Tao Jiang, Ming Li, and John Tromp. Lower bounds on learning decision lists and trees. Information and Computation, 126(2):114-122, 1996.
\bibitem{ref24} Johan Hastad. Clique is hard to approximate within n\textasciicircum{}1-$\varepsilon$. In Symposium on Foundations of Computer Science (FOCS), pages 627-636, 1996.
\bibitem{ref25} Shuichi Hirahara, Igor Carboni Oliveira, and Rahul Santhanam. NP-hardness of minimum circuit size problem for OR-AND-MOD circuits. In Computational Complexity Conference (CCC), pages 5:1-5:31, 2018.
\bibitem{ref26} Shuichi Hirahara and Osamu Watanabe. Limits of minimum circuit size problem as oracle. In Conference on Computational Complexity (CCC), pages 18:1-18:20, 2016.
\bibitem{ref27} John M. Hitchcock and Aduri Pavan. On the NP-completeness of the minimum circuit size problem. In Foundation of Software Technology and Theoretical Computer Science (FSTTCS), pages 236-245, 2015.
\bibitem{ref28} Rahul Ilango. Approaching MCSP from above and below: Hardness for a conditional variant and ac$^0$[p]. In Thomas Vidick, editor, Innovations in Theoretical Computer Science Conference (ITCS), volume 151 of LIPIcs, pages 34:1-34:26. Schloss Dagstuhl - Leibniz-Zentrum für Informatik, 2020.
\bibitem{ref29} Rahul Ilango. Constant depth formula and partial function versions of MCSP are hard. In Proceedings of the Symposium on Foundations of Computer Science (FOCS), pages 424-433. IEEE, 2020.
\bibitem{ref30} Rahul Ilango, Bruno Loff, and Igor Carboni Oliveira. Np-hardness of circuit minimization for multi-output functions. In Shubhangi Saraf, editor, Computational Complexity Conference (CCC), volume 169 of LIPIcs, pages 22:1-22:36. Schloss Dagstuhl - Leibniz-Zentrum für Informatik, 2020.
\bibitem{ref31} Russell Impagliazzo, Valentine Kabanets, and Ilya Volkovich. The Power of Natural Properties as Oracles. In Computational Complexity Conference (CCC), volume 102, pages 7:1-7:20, 2018.
\bibitem{ref32} J. Stephen Judd. Learning in networks is hard. In International Conference on Neural Networks (ICNN), volume 2, pages 685-692, 1987.
\bibitem{ref33} Valentine Kabanets and Jin-yi Cai. Circuit minimization problem. In Symposium on Theory of Computing (STOC), pages 73-79, 2000.
\bibitem{ref34} Mauricio Karchmer, Eyal Kushilevitz, and Noam Nisan. Fractional covers and communication complexity. SIAM Journal on Discrete Mathematics, 8(1):76-92, 1995.
\bibitem{ref35} Mauricio Karchmer and Avi Wigderson. Monotone circuits for connectivity require super-logarithmic depth. In Proceedings of the Twentieth Annual ACM Symposium on Theory of Computing, STOC '88, page 539–550. Association for Computing Machinery, 1988.
\bibitem{ref36} Subhash Khot and Rishi Saket. Hardness of minimizing and learning DNF expressions. In Symposium on Foundations of Computer Science (FOCS), pages 231-240, 2008.
\bibitem{ref37} Jan Krajícek. Forcing with Random Variables and Proof Complexity. Cambridge University Press, 2011.
\bibitem{ref38} Eyal Kushilevitz and Noam Nisan. Communication Complexity. Cambridge University Press, 1997.
\bibitem{ref39} Eyal Kushilevitz and Enav Weinreb. On the complexity of communication complexity. In Symposium on Theory of Computing (STOC), pages 465-474, 2009.
\bibitem{ref40} Yanyi Liu and Rafael Pass. On one-way functions and kolmogorov complexity. arXiv preprint, 2020. URL: http://arxiv.org/abs/2009.11514.
\bibitem{ref41} L. Lovasz and M. Saks. Lattices, mobius functions and communications complexity. In Symposium on Foundations of Computer Science (FOCS), SFCS '88, page 81–90, USA, 1988. IEEE Computer Society.
\bibitem{ref42} Carsten Lund and Mihalis Yannakakis. On the hardness of approximating minimization problems. Journal of the ACM, 41(5):960-981, 1994.
\bibitem{ref43} William J. Masek. Some NP-complete set covering problems. Unpublished Manuscript, 1979.
\bibitem{ref44} Dylan M. McKay, Cody D. Murray, and R. Ryan Williams. Weak lower bounds on resource-bounded compression imply strong separations of complexity classes. In Symposium on Theory of Computing (STOC), STOC 2019, page 1215–1225, New York, NY, USA, 2019. Association for Computing Machinery.
\bibitem{ref45} Dylan M. McKay, Cody D. Murray, and R. Ryan Williams. Weak lower bounds on resource-bounded compression imply strong separations of complexity classes. In Symposium on Theory of Computing (STOC), 2019.
\bibitem{ref46} Peter Bro Miltersen, Noam Nisan, Shmuel Safra, and Avi Wigderson. On data structures and asymmetric communication complexity. Journal of Computer and System Sciences, 57(1):37-49, 1998.
\bibitem{ref47} Moritz Müller and Ján Pich. Feasibly constructive proofs of succinct weak circuit lower bounds. Annals of Pure and Applied Logic, 171(2), 2020.
\bibitem{ref48} Cody D. Murray and Richard Ryan Williams. On the (non) NP-hardness of computing circuit complexity. In Conference on Computational Complexity (CCC), pages 365-380, 2015.
\bibitem{ref49} Joseph Naor and Moni Naor. Small-bias probability spaces: Efficient constructions and applications. SIAM journal on computing, 22(4):838-856, 1993.
\bibitem{ref50} Moni Naor and Omer Reingold. Number-theoretic constructions of efficient pseudo-random functions. Journal of the ACM, 51(2):231-262, 2004.
\bibitem{ref51} Moni Naor, Leonard J. Schulman, and Aravind Srinivasan. Splitters and Near-Optimal Derandomization. In Symposium on Foundations of Computer Science (FOCS), pages 182-191, 1995.
\bibitem{ref52} Noam Nisan. Pseudorandom bits for constant depth circuits. Combinatorica, 11(1):63-70, 1991.
\bibitem{ref53} Igor Carboni Oliveira, Ján Pich, and Rahul Santhanam. Hardness magnification near state-of-the-art lower bounds. In Conference on Computational Complexity (CCC), 2019.
\bibitem{ref54} Igor Carboni Oliveira and Rahul Santhanam. Conspiracies between learning algorithms, circuit lower bounds, and pseudorandomness. In Computational Complexity Conference (CCC), pages 18:1-18:49, 2017.
\bibitem{ref55} Igor Carboni Oliveira and Rahul Santhanam. Hardness magnification for natural problems. In Symposium on Foundations of Computer Science (FOCS), pages 65-76, 2018.
\bibitem{ref56} Denis Pankratov. Direct sum questions in classical communication complexity. Master’s thesis, University of Chicago, 2012.
\bibitem{ref57} Leonard Pitt and Leslie G. Valiant. Computational limitations on learning from examples. Journal of the ACM, 35(4):965-984, 1988.
\bibitem{ref58} Anup Rao and Amir Yehudayoff. Communication Complexity: and Applications. Cambridge University Press, 2020.
\bibitem{ref59} Alexander A Razborov. On submodular complexity measures. Boolean Function Complexity,(M. Paterson, Ed.), pages 76-83, 1992.
\bibitem{ref60} Alexander A Razborov and Steven Rudich. Natural proofs. Journal of Computer and System Sciences, 55(1):24-35, 1997.
\bibitem{ref61} Alexander A. Razborov and Steven Rudich. Natural proofs. Journal of Computer and System Sciences, 55(1):24-35, 1997.
\bibitem{ref62} Pranab Sen and Srinivasan Venkatesh. Lower bounds for predecessor searching in the cell probe model. Journal of Computer and System Sciences, 74(3):364-385, 2008.
\bibitem{ref63} Boris A Trakhtenbrot. A survey of Russian approaches to perebor (brute-force search) algorithms. Annals of the History of Computing, 6(4):384-400, 1984.
\bibitem{ref64} Christopher Umans, Tiziano Villa, and Alberto L. Sangiovanni-Vincentelli. Complexity of two-level logic minimization. IEEE Transactions on CAD of Integrated Circuits and Systems, 25(7):1230-1246, 2006.
\bibitem{ref65} David Zuckerman. Linear Degree Extractors and the Inapproximability of Max Clique and Chromatic Number. Theory of Computing, 3(1):103-128, 2007.
\end{thebibliography}
